\documentclass[11pt,oneside]{article}
\newcommand{\TopicTitle}{Apuntes del Tema 10: Redes Neuronales Recurrentes}
\usepackage[a4paper,margin=2.6cm,headheight=15pt]{geometry}
\usepackage{fontspec}
\usepackage{polyglossia}
\setdefaultlanguage{spanish}
\setotherlanguage{english}
\setmainfont{Latin Modern Roman}
\setsansfont{Latin Modern Sans}
\setmonofont{Latin Modern Mono}[Scale=MatchLowercase]
\usepackage{amsmath,amssymb,mathtools}
\usepackage{graphicx}
\usepackage{xcolor}
\usepackage{booktabs}
\usepackage{tabularx}
\usepackage{array}
\usepackage{longtable}
\usepackage{enumitem}
\usepackage{listings}
\usepackage{fancyhdr}
\usepackage{titlesec}
\usepackage{microtype}
\usepackage{hyperref}
\usepackage{bookmark}
\usepackage{caption}

\definecolor{ink}{HTML}{243447}
\definecolor{accent}{HTML}{176B87}
\definecolor{codebg}{HTML}{F3F5F7}
\definecolor{codegray}{HTML}{65727E}

\hypersetup{
  colorlinks=true,
  linkcolor=accent,
  urlcolor=accent,
  citecolor=accent,
  pdftitle={\TopicTitle},
  pdfauthor={Foundations of Machine Learning and Deep Learning}
}
\setcounter{secnumdepth}{3}
\setcounter{tocdepth}{3}
\setlist[itemize]{leftmargin=1.7em,itemsep=2pt,topsep=4pt}
\setlist[enumerate]{leftmargin=1.9em,itemsep=2pt,topsep=4pt}
\setlength{\parindent}{0pt}
\setlength{\parskip}{5pt plus 1pt minus 1pt}
\titleformat{\section}{\Large\sffamily\bfseries\color{ink}}{\thesection}{0.8em}{}
\titleformat{\subsection}{\large\sffamily\bfseries\color{accent}}{\thesubsection}{0.8em}{}
\titleformat{\subsubsection}{\normalsize\sffamily\bfseries\color{ink}}{\thesubsubsection}{0.8em}{}
\titlespacing*{\section}{0pt}{2.0em}{0.7em}
\titlespacing*{\subsection}{0pt}{1.4em}{0.4em}
\titlespacing*{\subsubsection}{0pt}{1.0em}{0.3em}
\lstset{
  backgroundcolor=\color{codebg},
  basicstyle=\ttfamily\small,
  breaklines=true,
  breakatwhitespace=false,
  columns=fullflexible,
  frame=single,
  framerule=0.3pt,
  rulecolor=\color{codegray},
  numbers=left,
  numbersep=8pt,
  numberstyle=\scriptsize\color{codegray},
  showstringspaces=false,
  keepspaces=true,
  tabsize=4,
  keywordstyle=\color{accent}\bfseries,
  commentstyle=\color{codegray}\itshape,
  stringstyle=\color{purple}
}
\captionsetup{font=small,labelfont=bf}
\pagestyle{fancy}
\fancyhf{}
\fancyhead[L]{\small\sffamily\color{codegray}Foundations of Machine Learning and Deep Learning}
\fancyhead[R]{\small\sffamily\color{codegray}\nouppercase{\rightmark}}
\fancyfoot[C]{\thepage}
\renewcommand{\headrulewidth}{0.4pt}
\renewcommand{\footrulewidth}{0pt}
\newcommand{\topiccover}[1]{%
  \begin{titlepage}
    \centering
    \vspace*{2.5cm}
    {\Large\sffamily\color{accent}APUNTES UNIVERSITARIOS\par}
    \vspace{1.5cm}
    {\Huge\sffamily\bfseries\color{ink}Foundations of Machine Learning\\[0.25em]
    and Deep Learning\par}
    \vspace{1.8cm}
    {\rule{0.72\textwidth}{1pt}\par}
    \vspace{1.2cm}
    {\LARGE\sffamily\bfseries #1\par}
    \vfill
    {\large\sffamily\color{codegray}\today\par}
  \end{titlepage}
}

\begin{document}
\topiccover{Apuntes del Tema 10: Redes Neuronales Recurrentes}
\tableofcontents
\clearpage
\markboth{Apuntes del Tema 10: Redes Neuronales Recurrentes}{Apuntes del Tema 10: Redes Neuronales Recurrentes}
\section{Apuntes del Tema 10: Redes Neuronales Recurrentes}

\subsection{1. De Datos Estáticos a Datos Secuenciales}

\subsubsection{1.1 Limitaciones de los Modelos Feedforward (MLPs y CNNs)}

\begin{itemize}

\item \textbf{Asunción de Muestras Independientes (\(i.i.d.\)):} Los modelos tradicionales como los Perceptrones Multicapa (MLP) o las Redes Convolucionales (CNN) asumen que las muestras de entrada son independientes e idénticamente distribuidas.

\end{itemize}

\begin{itemize}

\item \textbf{Entradas de Tamaño Fijo:} Un MLP o una CNN requieren un vector o tensor de dimensiones de entrada rígidamente predefinidas.

\end{itemize}

\begin{itemize}

\item \textbf{Incapacidad para Capturar Dependencias Temporales:} En datos con orden temporal o secuencial (como texto, audio, series temporales financieras o trayectorias GPS), la información del elemento actual depende de los elementos anteriores.

\end{itemize}

\subsubsection{1.2 Modelos Autorregresivos y el Estado Oculto}

Para modelar secuencias \(X = (x_1, x_2, \dots, x_T)\), deseamos estimar la probabilidad conjunta mediante la regla de la cadena de probabilidades:

\[P(x_1, x_2, \dots, x_T) = \prod_{t=1}^{T} P(x_t \mid x_1, \dots, x_{t-1})\]

\begin{itemize}

\item \textbf{Modelos de Marcov (Orden \(k\)):} Asumen que \(x_t\) solo depende de los \(k\) pasos anteriores.

\end{itemize}

\begin{itemize}

\item **Mecanismo de Estado Oculto (\emph{Hidden State}):\textbf{ En lugar de truncar la historia a \(k\) pasos, una Red Neuronal Recurrente (RNN) mantiene un }estado latente \(h_t\)** que actúa como una memoria comprimida de toda la secuencia procesada desde \(t=1\) hasta \(t-1\).

\end{itemize}

\subsection{2. Arquitectura de la Red Neuronal Recurrente Básica (RNN Vanilla)}

\subsubsection{2.1 Ecuaciones del Estado Oculto y Salida}

En cada paso de tiempo \(t\), la RNN recibe la entrada actual \(x_t \in \mathbb{R}^d\) y el estado oculto del paso anterior \(h_{t-1} \in \mathbb{R}^h\).

\begin{enumerate}

\item \textbf{Actualización del Estado Oculto:}

\end{enumerate}

\[h_t = \phi\left(x_t W_{xh} + h_{t-1} W_{hh} + b_h\right)\]

\begin{enumerate}

\item \textbf{Cálculo de la Salida:}

\end{enumerate}

\[\hat{y}_t = \text{Softmax}\left(h_t W_{hq} + b_q\right) \quad \text{o} \quad o_t = h_t W_{hq} + b_q\]

donde:

\begin{itemize}

\item \(W_{xh} \in \mathbb{R}^{d \times h}\): Matriz de pesos para la entrada.

\end{itemize}

\begin{itemize}

\item \(W_{hh} \in \mathbb{R}^{h \times h}\): Matriz de pesos para la transición recurrente de estados.

\end{itemize}

\begin{itemize}

\item \(W_{hq} \in \mathbb{R}^{h \times q}\): Matriz de pesos para la capa de salida.

\end{itemize}

\begin{itemize}

\item \(\phi(\cdot)\): Función de activación no lineal (típicamente \(\tanh\) o \(\text{ReLU}\)).

\end{itemize}

\begin{quote}**Compartición de Pesos (\emph{Weight Sharing} Temporal):** Las matrices de pesos \(\{W_{xh}, W_{hh}, W_{hq}\}\) son exactamente las mismas en todos los pasos de tiempo \(t=1, \dots, T\).\end{quote}

\begin{quote}\end{quote}

\begin{quote}\end{quote}

\begin{lstlisting}

Estructura Desplegada en el Tiempo (Unfolded RNN):



  y_1         y_2                 y_T

   ▲           ▲                   ▲

   │           │                   │

┌─────┐     ┌─────┐             ┌─────┐

│ h_1 │───► │ h_2 │───► ... ───►│ h_T │

└─────┘ W_hh└─────┘ W_hh        └─────┘

   ▲           ▲                   ▲

   │ W_xh      │ W_xh              │ W_xh

  x_1         x_2                 x_T



\end{lstlisting}

\subsection{3. Entrenamiento: Retropropagación a través del Tiempo (BPTT)}

\subsubsection{3.1 Backpropagation Through Time (BPTT)}

Para entrenar las RNNs, la red recurrente se "despliega" a lo largo de los \(T\) pasos temporales formando un grafo acíclico dirigido. La pérdida total es la suma de las pérdidas en cada paso de tiempo:

\[\mathcal{L}(W) = \sum_{t=1}^{T} \mathcal{L}_t(y_t, \hat{y}_t)\]

Para calcular el gradiente respecto a la matriz de transición recurrente \(W_{hh}\), se aplica la regla de la cadena acumulando las dependencias pasadas:

\[\frac{\partial \mathcal{L}_t}{\partial W_{hh}} = \sum_{k=1}^{t} \frac{\partial \mathcal{L}_t}{\partial h_t} \cdot \left( \prod_{j=k+1}^{t} \frac{\partial h_j}{\partial h_{j-1}} \right) \cdot \frac{\partial h_k}{\partial W_{hh}}\]

\subsubsection{3.2 Problema Crítico: Desvanecimiento y Explosión del Gradiente}

El término \(\prod_{j=k+1}^{t} \frac{\partial h_j}{\partial h_{j-1}}\) implica multiplicar repetidamente la matriz de pesos \(W_{hh}^\top\).

\begin{itemize}

\item **Desvanecimiento del Gradiente (\emph{Vanishing Gradients}):\textbf{ Si los autovalores de \(W_{hh}\) son menores a \(1\), la magnitud del gradiente decrece exponencialmente hacia cero cuando la secuencia es larga (\(T \gg 0\)). La red }olvida las dependencias a largo plazo**.

\end{itemize}

\begin{itemize}

\item **Explosión del Gradiente (\emph{Exploding Gradients}):\textbf{ Si los autovalores son mayores a \(1\), el gradiente crece exponencialmente destruyendo el entrenamiento. Se soluciona mediante }Gradient Clipping**:

\end{itemize}

\[g \leftarrow \min\left(1, \frac{\theta}{\Vert{}g\Vert{}}\right) g\]

\subsection{4. Redes Gated: LSTM y GRU}

Para solucionar el desvanecimiento de gradiente y permitir recordar información a largo plazo, se crearon arquitecturas basadas en puertas de control (\emph{gates}).

\subsubsection{4.1 LSTM (\emph{Long Short-Term Memory})}

La LSTM sustituye la celda estándar por una estructura de memoria que mantiene un \textbf{estado de celda \(C_t\)} (\emph{Cell State} o autopista de información) y un \textbf{estado oculto \(h_t\)}.

\begin{lstlisting}

             ┌──────────────────────────────────────────────┐ (Cell State C_t)

             │                 x (f_t)        + (i_t * C~_t)│

C_{t-1} ─────┴─────────────────┴──────────────┴─────────────┴────────► C_t

                   ▲                ▲         ▲          ▲

                   │                │         │          │

h_{t-1} ─┬─────────┼───────────────┼─────────┼──────────┼─────────┬──► h_t

         │         │                │         │          │         │

x_t ─────┴─► [Forget Gate]    [Input Gate] [Candidate] [Output Gate]



\end{lstlisting}

\subsubsection{Puertas y Ecuaciones Matemáticas de la LSTM:}

\begin{enumerate}

\item \textbf{Forget Gate (\(f_t\)):} Decide qué información descartar del estado de celda anterior.

\end{enumerate}

\[f_t = \sigma\left(x_t W_{xf} + h_{t-1} W_{hf} + b_f\right)\]

\begin{enumerate}

\item \textbf{Input Gate (\(i_t\)) y Candidato (\(\tilde{C}_t\)):} Deciden qué nueva información almacenar.

\end{enumerate}

\[i_t = \sigma\left(x_t W_{xi} + h_{t-1} W_{hi} + b_i\right)\]

\[\tilde{C}_t = \tanh\left(x_t W_{xc} + h_{t-1} W_{hc} + b_c\right)\]

\begin{enumerate}

\item \textbf{Actualización del Cell State (\(C_t\)):}

\end{enumerate}

\[C_t = f_t \odot C_{t-1} + i_t \odot \tilde{C}_t\]

\begin{enumerate}

\item \textbf{Output Gate (\(o_t\)) y Estado Oculto (\(h_t\)):}

\end{enumerate}

\[o_t = \sigma\left(x_t W_{xo} + h_{t-1} W_{ho} + b_o\right)\]

\[h_t = o_t \odot \tanh(C_t)\]

\subsubsection{4.2 GRU (\emph{Gated Recurrent Unit})}

Variante simplificada de la LSTM que combina el estado de celda y el estado oculto en una sola variable \(h_t\), reduciendo el coste computacional y el número de parámetros.

\begin{enumerate}

\item \textbf{Reset Gate (\(r_t\)):} Determina cuánta información del pasado olvidar:

\end{enumerate}

\[r_t = \sigma\left(x_t W_{xr} + h_{t-1} W_{hr} + b_r\right)\]

\begin{enumerate}

\item \textbf{Update Gate (\(z_t\)):} Actúa simultáneamente como puerta de olvida e entrada:

\end{enumerate}

\[z_t = \sigma\left(x_t W_{xz} + h_{t-1} W_{hz} + b_z\right)\]

\begin{enumerate}

\item \textbf{Candidato de Estado Oculto (\(\tilde{h}_t\)):}

\end{enumerate}

\[\tilde{h}_t = \tanh\left(x_t W_{xh} + (r_t \odot h_{t-1}) W_{hh} + b_h\right)\]

\begin{enumerate}

\item \textbf{Estado Oculto Final (\(h_t\)):}

\end{enumerate}

\[h_t = (1 - z_t) \odot h_{t-1} + z_t \odot \tilde{h}_t\]

\subsection{5. RNNs Bidireccionales y Apiladas (\emph{Deep RNNs})}

\begin{itemize}

\item \textbf{RNN Bidireccional (BiRNN / BiLSTM):} Procesa la secuencia en ambas direcciones simultáneamente usando dos estados ocultos independientes (\(\overrightarrow{h}_t\) hacia adelante y \(\overleftarrow{h}_t\) hacia atrás). El estado combinado es \(h_t = [\overrightarrow{h}_t \,;\, \overleftarrow{h}_t]\).

\end{itemize}

\begin{itemize}

\item \emph{Uso:} Crucial en tareas de procesamiento de lenguaje natural (NLP) donde el contexto futuro influye en el significado del presente.

\end{itemize}

\begin{itemize}

\item **RNN Profunda / Apilada (\emph{Stacked RNN}):** Apila múltiples capas recurrentes donde la salida del estado oculto de la capa \(l\) sirve como entrada \(x_t\) para la capa \(l+1\).

\end{itemize}

\subsection{6. Práctica y Aplicación en Código (PyTorch)}

A continuación se incluye una implementación modular completa que ilustra la construcción de una red \textbf{LSTM Bidireccional} para la clasificación de secuencias (ej. análisis de sentimiento o series temporales) en PyTorch:

\begin{lstlisting}[language=Python]

import torch

import torch.nn as nn

import torch.optim as optim



# --------------------------------------------------------------------------

# 1. Definición del Modelo BiLSTM / GRU en PyTorch

# --------------------------------------------------------------------------

class SequenceClassifier(nn.Module):

    def __init__(self, vocab_size=1000, embed_dim=64, hidden_dim=128, num_classes=2, num_layers=2, dropout=0.3):

        super(SequenceClassifier, self).__init__()

        

        # Capa de Embedding para convertir índices de tokens en vectores continuos

        self.embedding = nn.Embedding(vocab_size, embed_dim)

        

        # Capa LSTM Bidireccional de múltiples capas

        self.lstm = nn.LSTM(

            input_size=embed_dim,

            hidden_size=hidden_dim,

            num_layers=num_layers,

            batch_first=True,       # Entrada esperada: (Batch, Seq_Len, Embed_Dim)

            bidirectional=True,     # BiLSTM

            dropout=dropout if num_layers > 1 else 0.0

        )

        

        # Capa lineal final: hidden_dim * 2 por la bidireccionalidad

        self.fc = nn.Linear(hidden_dim * 2, num_classes)

        self.dropout = nn.Dropout(p=dropout)



    def forward(self, x):

        # x shape: (Batch, Seq_Len)

        embedded = self.embedding(x) # (Batch, Seq_Len, Embed_Dim)

        

        # output shape: (Batch, Seq_Len, Hidden_Dim * 2)

        # h_n, c_n shape: (Num_Layers * 2, Batch, Hidden_Dim)

        lstm_out, (h_n, c_n) = self.lstm(embedded)

        

        # Concatenar el último estado oculto forward y backward de la última capa

        hidden_forward = h_n[-2, :, :]

        hidden_backward = h_n[-1, :, :]

        last_hidden = torch.cat((hidden_forward, hidden_backward), dim=1) # (Batch, Hidden_Dim * 2)

        

        out = self.dropout(last_hidden)

        logits = self.fc(out) # (Batch, Num_Classes)

        return logits



# --------------------------------------------------------------------------

# 2. Ejemplo de Bucle de Entrenamiento

# --------------------------------------------------------------------------

def train_rnn_demo():

    device = torch.device("cuda" if torch.cuda.is_available() else "cpu")

    

    # Hiperparámetros

    VOCAB_SIZE = 1000

    BATCH_SIZE = 16

    SEQ_LEN = 32

    NUM_CLASSES = 2

    

    # Modelo, pérdida y optimizador

    model = SequenceClassifier(

        vocab_size=VOCAB_SIZE, 

        embed_dim=64, 

        hidden_dim=128, 

        num_classes=NUM_CLASSES, 

        num_layers=2

    ).to(device)

    

    criterion = nn.CrossEntropyLoss()

    optimizer = optim.Adam(model.parameters(), lr=0.001)

    

    # Datos sintéticos de secuencia: (Batch, Seq_Len)

    dummy_input = torch.randint(0, VOCAB_SIZE, (BATCH_SIZE, SEQ_LEN)).to(device)

    dummy_labels = torch.randint(0, NUM_CLASSES, (BATCH_SIZE,)).to(device)

    

    model.train()

    optimizer.zero_grad()

    

    # Forward Pass

    logits = model(dummy_input)

    loss = criterion(logits, dummy_labels)

    

    # Backward Pass

    loss.backward()

    

    # Gradient Clipping esencial para evitar la explosión de gradiente en RNNs

    nn.utils.clip_grad_norm_(model.parameters(), max_norm=5.0)

    

    optimizer.step()

    

    print(f"Demostración RNN completada en dispositivo: {device}")

    print(f"Dimensión de salida (Logits): {logits.shape}")

    print(f"Pérdida inicial (Loss): {loss.item():.4f}")



if __name__ == "__main__":

    train_rnn_demo()



\end{lstlisting}

Dime si necesitas que profundicemos en algún aspecto de las \textbf{RNNs, LSTMs o GRUs} o si te gustaría preparar un cuestionario interactivo para poner a prueba estos conceptos de cara a tu máster.
\end{document}
